% §18 Exercises
%
% Cross-reference resolution status:
%
% DEFINED here:
%   \label{ex:ch18-1} through \label{ex:ch18-11}
%
% RESOLVED (used here, defined elsewhere):
%   \Cref{def:per-document-bloom}             -- §18.2
%   \Cref{def:inverted-index-bloom}           -- §18.2
%   \Cref{sec:bloom-based}                    -- §18.2
%   \Cref{sec:confidentiality-measures}       -- §18.3
%   \Cref{sec:adversarial-estimation}         -- §18.4
%   \Cref{sec:toy-system}                     -- §18.5
%   \Cref{prop:multi-query-leakage}           -- §18.3
%   \Cref{def:ppv}                            -- chapter 7 PPV
%   \Cref{def:f1}                             -- chapter 7 F1
%   \Cref{sec:derived-measures}               -- chapter 7 §7.4
%   \Cref{sec:confidence-intervals}           -- chapter 7 §7.3
%   \Cref{thm:chernoff}                       -- chapter 11 §11.4
%   \Cref{thm:bloom-optimal-k}                -- chapter 5 §5.2
%   \Cref{ax:element-indep}                   -- chapter 4 Axiom 1
%   \Cref{sec:maximizing-confidentiality}     -- chapter 17 §17.6
%   \Cref{def:cipher-map}                     -- chapter 17 §17.1
%   \Cref{ch:perfect-hash-filters}            -- chapter 15
%
% 11 exercises, 3 star / 5 starstar / 3 starstarstar:
% - star x 3: per-document Bloom space; confidentiality at single
%   workload; query-token computation in two sentences
% - starstar x 5: notebook-tied precision verification; per-doc vs
%   inverted-index space and query time; multi-query leakage bound;
%   adversarial defender iid Chernoff vs MAB UCB; F1 leakage on
%   balanced workload
% - starstarstar x 3: tuning fprate for target confidentiality;
%   per-doc vs BPHF storage; design rule synthesizing chapters 6,
%   7, 8, 14, 17, 18

% Difficulty-marker macros (matching chapters 1-17):
\providecommand{\exstar}{$\star$}
\providecommand{\exstarstar}{$\star\star$}
\providecommand{\exstarstarstar}{$\star\star\star$}

\section*{Exercises}
\addcontentsline{toc}{section}{Exercises}

\begin{enumerate}

% -----------------------------------------------------------------------
% WARM-UP  (star x 3): exercises 1, 2, 3
% -----------------------------------------------------------------------

\item \exstar
\label{ex:ch18-1}
\textbf{Per-document Bloom space.}
A corpus has $|D| = 10{,}000$ documents with $W_d = 50$ keywords
per document on average, and the per-document Bloom-filter
encrypted-search construction (\Cref{def:per-document-bloom}) is
configured at $\fprate = 0.01$ with the framework's
$\fnrate = 0$.
\begin{enumerate}
  \item[(a)] Apply the chapter~5 optimal-$k$ formula
    (\Cref{thm:bloom-optimal-k}) to compute $m_d$, the
    per-document filter size in bits.
  \item[(b)] Compute the total storage cost of the encrypted
    index in bits and in MB.
  \item[(c)] State in one sentence whether this is competitive
    with a plaintext keyword index that stores each keyword as
    8 bytes.
\end{enumerate}

\item \exstar
\label{ex:ch18-2}
\textbf{Confidentiality at a single workload.}
A per-document Bloom-filter encrypted-search system runs at
$\fprate = 0.05, \fnrate = 0$, and the workload has prevalence
$\pi = 0.1$ (one queried keyword in ten truly appears in the
queried document).
\begin{enumerate}
  \item[(a)] Substitute into the chapter~7 PPV formula
    (\Cref{def:ppv}) to compute $\PPV$ numerically.
  \item[(b)] Compute the confidentiality $C = 1 - \PPV$ from
    \Cref{sec:maximizing-confidentiality}.
  \item[(c)] State in one sentence what $C$ means operationally
    at this operating point.
\end{enumerate}

\item \exstar
\label{ex:ch18-3}
\textbf{Query-token computation.}
In the toy system of \Cref{sec:toy-system}, the client computes a
query token before sending it to the server. State in two
sentences (a) what the client computes as the query token (be
specific about the cryptographic function used) and (b) what
property of that function ensures the server cannot recover the
plaintext keyword from the token.

% -----------------------------------------------------------------------
% MEDIUM  (starstar x 5): exercises 4, 5, 6, 7, 8
% -----------------------------------------------------------------------

\item \exstarstar
\label{ex:ch18-4}
\textbf{Notebook-tied precision verification.}
Use the chapter~18 notebook,
\texttt{python/notebooks/ch18-encrypted-search.ipynb}, on a
synthetic corpus of $|D| = 1000$ documents at $\fprate = 0.01$,
$\fnrate = 0$. Query a keyword that truly appears in $100$
documents.
\begin{enumerate}
  \item[(a)] Run the query and compute the empirical precision
    (true positives over total reported matches).
  \item[(b)] Compare with the chapter~7 PPV prediction at
    $\pi = 0.1$ from \Cref{sec:confidentiality-measures}.
    Report the absolute difference.
  \item[(c)] Repeat over $100$ independent corpora and average.
    Comment on convergence to the PPV prediction.
\end{enumerate}

\item \exstarstar
\label{ex:ch18-5}
\textbf{Per-document versus inverted-index.}
Take a corpus of $|D| = 10^6$ documents with $50$ keywords each
and a vocabulary of $|W| = 10^5$ keywords. Both Bloom
constructions of \Cref{sec:bloom-based} run at
$\fprate = 0.01$, $\fnrate = 0$.
\begin{enumerate}
  \item[(a)] Compute the per-document construction's total
    storage cost in bits using the chapter~5 optimal-$k$
    formula.
  \item[(b)] Compute the inverted-index construction's total
    storage cost (\Cref{def:inverted-index-bloom}), using
    $|D_w| \approx |D| \cdot W_d / |W|$.
  \item[(c)] Compare the two server-side query-time costs and
    state which scales better for high-throughput workloads.
\end{enumerate}

\item \exstarstar
\label{ex:ch18-6}
\textbf{Multi-query leakage bound.}
A user issues $m = 1000$ queries against an encrypted-search
system at $\fprate = 0.01$, $\fnrate = 0$, with the queries
drawn from a fixed distribution (the chapter~4 independence
axiom \Cref{ax:element-indep} is in force).
\begin{enumerate}
  \item[(a)] Apply the chapter~18 multi-query leakage scaling
    proposition (\Cref{prop:multi-query-leakage}) to bound the
    server's posterior precision on any single $(w, d_i)$
    plaintext pair after the batch.
  \item[(b)] At workload prevalence $\pi = 0.05$, compute the
    per-query PPV. State the rate at which the
    confidentiality bits per query decay across the batch.
  \item[(c)] Identify the two practical mitigations the
    multi-query leakage analysis suggests: a fixed query
    budget or a key-rotation schedule. State in one sentence
    which mitigation is easier to deploy operationally.
\end{enumerate}

\item \exstarstar
\label{ex:ch18-7}
\textbf{Adversarial defender: iid Chernoff versus MAB UCB.}
A defender monitors a deployed system at $100$ queries per
minute for one hour, observing $540$ matches across $6000$
queries. The vocabulary has $K = 1000$ keywords actively
queried.
\begin{enumerate}
  \item[(a)] Compute the iid Wilson CI at $95\%$ confidence
    (\Cref{sec:confidence-intervals}). Substitute the upper
    end into the confidentiality formula at $\fprate = 0.01$,
    $\fnrate = 0$ and report the iid lower bound on $C$.
  \item[(b)] Compute the MAB UCB upper bound on the worst-arm
    rate at joint confidence $1 - \delta = 0.95$ across $K =
    1000$ arms (union-bounded). Report the MAB lower bound on
    $C$.
  \item[(c)] Identify which bound is correct and state in one
    sentence why the other fails on adversarial traffic.
\end{enumerate}

\item \exstarstar
\label{ex:ch18-8}
\textbf{$\Fone$ leakage on a balanced workload.}
A balanced workload has $\pi = 0.5$ (half the queried keywords
truly appear in the queried document). The encrypted-search
system runs at $\fprate = 0.01$, $\fnrate = 0$.
\begin{enumerate}
  \item[(a)] Compute $\PPV$ at this operating point.
  \item[(b)] Apply the chapter~7 $\Fone$ formula
    (\Cref{def:f1}) to compute the system's $\Fone$-score
    at $\fnrate = 0$ (recall is $1$, precision is $\PPV$).
  \item[(c)] State in one sentence how the $\Fone$-leakage
    bound from \Cref{sec:derived-measures} differs from the
    PPV bound when the workload is balanced rather than
    rare-positive.
\end{enumerate}

% -----------------------------------------------------------------------
% CHALLENGING  (starstarstar x 3): exercises 9, 10, 11
% -----------------------------------------------------------------------

\item \exstarstarstar
\label{ex:ch18-9}
\textbf{Tuning $\fprate$ for a confidentiality target.}
A deployment requires $C \ge 0.9$ at prevalence $\pi = 0.05$.
The cipher map has $\fnrate = 0$.
\begin{enumerate}
  \item[(a)] Solve $\fprate^* = \pi C / [(1 - \pi)(1 - C)]$
    for the minimum $\fprate$ meeting the target. Report the
    value.
  \item[(b)] Compute the chapter~5 filter size $m_d$ at
    $W_d = 10$. Compare with the size at $\fprate = 0.01$.
  \item[(c)] Quantify the client's verification load as the
    expected number of false-positive matches per query
    against a $|D| = 10^5$-document corpus.
\end{enumerate}

\item \exstarstarstar
\label{ex:ch18-10}
\textbf{Per-document Bloom versus BPHF inverted index.}
The chapter~15 BPHF gives an FPR-zero construction for
inverted-index lookups when the keyword-to-document mapping is
deterministic and the query distribution is known in advance.
Take the corpus from Exercise~5.
\begin{enumerate}
  \item[(a)] Estimate the BPHF storage cost using the
    chapter~15 $O(|W| \cdot \log |D|)$ space bound for a
    minimal perfect hash with rank information.
  \item[(b)] Compare with the per-document Bloom storage cost
    from Exercise~5(a) at $\fprate = 0.01$.
  \item[(c)] State the cross-over $\fprate$ at which the two
    storage costs are equal. (See
    \Cref{ch:perfect-hash-filters} for the BPHF cost model.)
\end{enumerate}

\item \exstarstarstar
\label{ex:ch18-11}
\textbf{The framework as design language.}
State a design rule for choosing $\fprate$, $\fnrate$, and the
operating prevalence $\pi$ in an encrypted-search deployment,
given a privacy budget, a correctness budget, and a storage
budget.
\begin{enumerate}
  \item[(a)] Write the rule as a function from the three
    budgets to $(\fprate^*, \fnrate^*, \pi^*)$. The codomain
    may include ``no feasible operating point exists.''
  \item[(b)] Justify the rule via the book's chapters: which
    chapter prices each budget, and which equation the rule
    is solving. Reference at least chapters~6, 7, 8, 14, 17,
    and~18.
  \item[(c)] Apply the rule to a deployment with $C \ge 0.5$,
    expected false-positive matches per query at most $10$
    over a $10^5$-document corpus, and storage budget $1$ MB.
    Report the operating point or the infeasibility
    determination.
\end{enumerate}

\end{enumerate}
